% TOP_PROBLEM: {"id": 30007281, "title": "Constructible tame eigensheaves with regular monodromy", "category_id": 1, "category": {"id": 1, "name": "number_theory", "display_name": "Number Theory", "description": "Properties of integers, prime numbers, Diophantine equations.", "slug": "number-theory", "order_index": 1, "created_at": "2026-07-31T15:26:25.670Z"}, "status": "open"}
% FIRST_POSTED: 2026-10-06
% ATTEMPT_STATUS: solved
% OPENAI_CLAIM: {"schema_version":1,"source_repo":"https://github.com/openai/math","source_commit":"adc7f1241b42e322a6451854ab7e4b4c146bf78a","release_date":"2026-10-06","match":"full","resolution":"proved","independently_reviewed":false,"second_pass":"agreed","families":[{"family":"014","title":"Restricted geometric Langlands, global Arthur enhancements, and generic Ramanujan","subject":"Number theory","manuscripts":[{"title":"Constructible tame Hecke eigensheaves in positive characteristic","pdf_path":"preprints/Constructible-tame-Hecke-eigensheaves-in-positive-characteristic-October-5-2026/constructible-tame-hecke-eigensheaves-positive-characteristic.pdf","readme_path":"preprints/Constructible-tame-Hecke-eigensheaves-in-positive-characteristic-October-5-2026/README.md","date":"2026-10-05","theorem_ref":"preprints/Constructible-tame-Hecke-eigensheaves-in-positive-characteristic-October-5-2026/build/manuscript/introduction.tex :: thm:main"}],"lean":null,"reasoning_trace":null}]}
\documentclass[11pt]{article}
\usepackage[margin=25mm]{geometry}
\usepackage{fontspec}
\setmainfont{Latin Modern Roman}
\usepackage{amsmath,amssymb,mathtools}
\usepackage{unicode-math}
\setmathfont{Latin Modern Math}
\usepackage{enumitem}
\usepackage{xurl,hyperref,bookmark}
\hypersetup{colorlinks=true,urlcolor=blue}
\setcounter{secnumdepth}{0}
\setlength{\parindent}{0pt}
\setlength{\parskip}{5pt}
\setlength{\emergencystretch}{3em}
\begin{document}
\section*{Constructible tame eigensheaves with regular monodromy}

\subsection{OpenAI's claim}
OpenAI presents the following claims relevant to this question, with manuscript numbers following paper order in Sources. In manuscript 1 (family 014; 2026-10-05), the asserted result is that tame regular-unipotent, geometrically dense parameters admit nonzero constructible perverse Hecke eigensheaves with one Borel-level point for the specified simply connected groups and characteristic hypotheses.

\subsection{Outline of the argument}
For manuscript 1 (family 014; 2026-10-05), the parameter is glued through a separating node to an unramified family. Two specializations and transverse Hecke cuts detect a nonzero flag-level object; unipotent monodromy on the enhancement torus permits descent to ordinary Borel level.

\subsection{Scope relative to this problem}
The result supplies the conclusion in this problem with its displayed hypotheses and quantifiers. The broader statements and applications in the manuscript are kept separate from the particular question above. The supporting result is manuscript 1 (family 014; 2026-10-05).

\subsection{Formal verification}
No Lean formalization is listed at this snapshot for manuscript 1 (family 014; 2026-10-05).

This site has not executed the Lean code. Described statement coverage is not an independent proof check.

\subsection{Status caveat}
The manuscripts were produced by an internal OpenAI model. OpenAI cautions that results without formalization can contain mistakes. This site has not independently checked the proofs; the label records the scope of OpenAI's claim.

\subsection{Sources}
\begin{itemize}
\item OpenAI, \emph{Constructible tame Hecke eigensheaves in positive characteristic}, 2026-10-05; citation key \texttt{OAI:Constructible-tame-Hecke-eigensheaves-in-positive-characteristic-October-5-2026}. \href{https://github.com/openai/math/blob/adc7f1241b42e322a6451854ab7e4b4c146bf78a/preprints/Constructible-tame-Hecke-eigensheaves-in-positive-characteristic-October-5-2026/constructible-tame-hecke-eigensheaves-positive-characteristic.pdf}{PDF (GitHub)}; \href{https://github.com/openai/math/raw/adc7f1241b42e322a6451854ab7e4b4c146bf78a/preprints/Constructible-tame-Hecke-eigensheaves-in-positive-characteristic-October-5-2026/constructible-tame-hecke-eigensheaves-positive-characteristic.pdf}{PDF (direct)}; \href{https://github.com/openai/math/blob/adc7f1241b42e322a6451854ab7e4b4c146bf78a/preprints/Constructible-tame-Hecke-eigensheaves-in-positive-characteristic-October-5-2026/README.md}{Manuscript page and citation}; \href{https://github.com/openai/math/blob/main/preprints/Constructible-tame-Hecke-eigensheaves-in-positive-characteristic-October-5-2026/constructible-tame-hecke-eigensheaves-positive-characteristic.pdf}{Latest version}.
\item \href{https://github.com/openai/math/blob/adc7f1241b42e322a6451854ab7e4b4c146bf78a/CONTENTS.md}{OpenAI catalogue entry}.
\end{itemize}
\end{document}
